% Appendix: an illustrative note on Andrewes' sources. DRAFT — Wilson's to edit.
% GROUNDING. Every page reference here was checked against our own transcripts,
% and every attribution is either (a) printed in the 1853 plate, (b) supplied by
% the 1853 editor in square brackets, or (c) identified by Brightman 1903 and
% recorded in apparatus/BRIGHTMAN-collation.md. The distinction is load-bearing
% and the text keeps it. The count of scripture references is from the plate.
% Nothing here is an attribution of our own making.
\begingroup
\setlength{\parskip}{0.55\baselineskip}
\setlength{\parindent}{0pt}
\addcontentsline{toc}{section}{Appendix: Some of the Sources}
{\centering\Large\scshape Appendix\par}
\vspace{0.4\baselineskip}
{\centering\itshape Some of the sources Andrewes drew on\par}
\vspace{1.5\baselineskip}

This note is illustrative and not a source apparatus. It names enough of what
Andrewes gathered to show the kind of book he was making, and stops well short of
what a full accounting would require.

Three sorts of attribution are mixed in what follows, and they are not of equal
weight. Some authorities are printed in the 1853 itself, standing in the margin at
the head of the sentence they belong to --- \emph{Cyprian.}, \emph{Beda.},
\emph{Seneca.} Some are the 1853 editor's own identifications, which he set in
square brackets, and which are printed here in brackets for that reason. And some
were identified neither by Andrewes nor by his editor but by F.\,E. Brightman in
1903, working from the sources; where a line below rests on Brightman it says so.
The page numbers are the 1853's, as printed on the inner edge of every leaf here.

\vspace{0.8\baselineskip}
{\scshape Scripture, before anything else.}\quad
The 1853 prints better than two thousand scripture references. That number is the
plainest fact about this book: it is a cento, and the greater part of what
Andrewes prays is quotation. The text he quotes is the Septuagint and the Vulgate,
not any English Bible --- which is why the translation here reaches for
Coverdale's Psalter ahead of the Authorised Version, as the preface explains. The
Prayer of Manasses is drawn on at printed 171 (the editor's
identification, [\emph{Orat.\ Manassis.}]).

\vspace{0.8\baselineskip}
{\scshape The liturgies.}\quad
The Liturgy of S.\ Chrysostom is named twice in Part I, at printed 33 and 55, in
the plate itself. The \emph{Sursum corda} that opens the epigraph leaf of Part III
at printed 398 is older than any of the liturgies Andrewes could have held: it
stands in the Hippolytean canons and in Cyprian's treatise on the Lord's Prayer,
an identification we owe to Brightman.

\vspace{0.8\baselineskip}
{\scshape The creeds and the Prayer Book.}\quad
The Nicene Creed is named in the plate at printed 177. The \emph{Te Deum} is drawn
on at printed 25 and again at 431, on the editor's identification. And at printed
199 the margin reads [\emph{Litan.\ Angl.}] --- the English Litany. A bishop's
private book, kept in Greek and Latin, quietly quoting the vernacular liturgy he
had to lead on Sundays.

\vspace{0.8\baselineskip}
{\scshape The Greek fathers.}\quad
Chrysostom again at printed 417. Theophylact opens Part III at printed 398, and
Brightman located the passage exactly: the commentary on S.\ Luke xviii, where the
Greek \emph{βάθρον καὶ κρηπίς} is rendered by the Latin \emph{fundamentum et
basis} --- foundation and ground --- which is the first line of that leaf. Gregory
of Nyssa stands in the catena at printed 386, from the first homily on the Lord's
Prayer.

\vspace{0.8\baselineskip}
{\scshape The Latin fathers.}\quad
Part II gathers them in a short space, each named in the plate at the head of its
sentence: Cyprian at printed 381, then Bede, Arnobius and Jerome at 382, then
Augustine twice at 383 --- once from \emph{contra Maximinum} and once from the
first book of the \emph{Confessions} --- and Augustine again at 417 and 422.

One of these is worth a note of its own. At printed 380 the plate attributes a
passage \emph{ex Fulgentii l.\ i.\ ad Mon.}, to the first book of Fulgentius of
Ruspe \emph{to Monimus}. Our own reading suspected the passage was Augustine's,
and the attribution was kept as printed rather than mended. Brightman, working
independently from the sources, gives Fulgentius too. The plate was right and the
doubt was ours.

\vspace{0.8\baselineskip}
{\scshape A medieval archbishop.}\quad
Printed 370 opens \emph{Oratio Tho.\ Bradwardini} --- Thomas Bradwardine,
archbishop of Canterbury for five weeks in 1349. Brightman identifies the work as
\emph{de causa Dei contra Pelagium}, and the edition Andrewes is likely to have
used as Savile's, London 1618.

\vspace{0.8\baselineskip}
{\scshape Pagan antiquity, and no apology for it.}\quad
Cicero is named in the plate at printed 385 and again at 386, the second time for
Cato, who called himself to account each night for the day's business ---
\emph{de senectute}, on Brightman's identification. Seneca is named at printed
385. At 386 Andrewes cites \emph{Ausonius out of Pythagoras}, a quotation at one
remove and marked as such.

Two things about this stretch repay attention. The first is that nothing in the
page ranks these authorities below the Fathers or above them. Bede, Augustine,
Cicero and Seneca stand in one column, each with his sentence, and the reader is
left to make of that what he will. The second is smaller and more human: the
sentence Seneca gives --- that you must catch yourself out before you can mend
yourself --- Seneca himself had from Epicurus. The chain is Epicurus to Seneca to
Andrewes, and a reader coming on that page would very likely take the thought for
a Father's.

\vspace{0.8\baselineskip}
{\scshape Jewish and contemporary.}\quad
At printed 386, in the middle of his own sentence rather than in a margin,
Andrewes writes \emph{Recte igitur Rab.\ J.} --- rightly, then, Rabbi J. --- on
not putting repentance off until tomorrow. Hebrew stands in his text in more
places than the 1853 prints; what it left out is recorded in the apparatus, and
the preface says why it has not been restored here. On the following page,
printed 387, he cites \emph{Rev.\ Usserius} --- James Ussher, then a
contemporary --- with page numbers, as one cites a book lately read.

\vspace{0.8\baselineskip}
{\scshape A caution about the margins.}\quad
The attributions printed in the plate are not always exact, and they have been
left as they are. The clearest instance is the \emph{Seneca.} at printed 385,
which stands at the head of two sentences. The first is Seneca's, from the
twenty-eighth moral letter. The second --- that a sore never found festers and
goes uncured --- is not in that letter and has not been located anywhere else.
It may be Andrewes' own amplification of Seneca's metaphor; it may equally come
from a collection of sentences of the kind he plainly used. The honest statement
is that it is unattested elsewhere, and that one margin covers two sentences of
which only the first is certainly Seneca's. A margin in this book is a pointer,
not a warranty.

\endgroup
\clearpage
